% !TEX program = xelatex
% 기존 weekly_presentation_template_4x3.tex 형식 사용.
% docs 폴더에서 XeLaTeX로 두 번 컴파일합니다.
\documentclass[aspectratio=43,11pt,t]{beamer}
\usepackage{fontspec}
\usepackage{kotex}
\usepackage{graphicx}
\usepackage{tikz}
\usetikzlibrary{arrows.meta}
\usepackage{array}
\usepackage{booktabs}
\usepackage{fancyvrb}
\setmainfont{D2Coding-Regular.ttf}[Path=assets/lecture/,BoldFont={D2Coding-Bold.ttf},ItalicFont={D2Coding-Regular.ttf},BoldItalicFont={D2Coding-Bold.ttf}]
\setmainhangulfont{D2Coding-Regular.ttf}[Path=assets/lecture/,BoldFont={D2Coding-Bold.ttf},ItalicFont={D2Coding-Regular.ttf},BoldItalicFont={D2Coding-Bold.ttf}]
\setsansfont{D2Coding-Regular.ttf}[Path=assets/lecture/,BoldFont={D2Coding-Bold.ttf},ItalicFont={D2Coding-Regular.ttf},BoldItalicFont={D2Coding-Bold.ttf}]
\setsanshangulfont{D2Coding-Regular.ttf}[Path=assets/lecture/,BoldFont={D2Coding-Bold.ttf},ItalicFont={D2Coding-Regular.ttf},BoldItalicFont={D2Coding-Bold.ttf}]
\setmonofont{D2Coding-Regular.ttf}[Path=assets/lecture/,BoldFont={D2Coding-Bold.ttf},ItalicFont={D2Coding-Regular.ttf},BoldItalicFont={D2Coding-Bold.ttf}]
\definecolor{UOSBlue}{HTML}{005EB8}
\definecolor{TextGray}{HTML}{303030}
\setbeamertemplate{navigation symbols}{}
\setbeamercolor{normal text}{fg=TextGray,bg=white}
\setbeamercolor{frametitle}{fg=UOSBlue,bg=white}
\setbeamercolor{structure}{fg=UOSBlue}
\setbeamercolor{block title}{fg=white,bg=UOSBlue}
\setbeamercolor{block body}{fg=TextGray,bg=UOSBlue!6}
\setbeamertemplate{itemize item}{\color{UOSBlue}\small$\bullet$}
\setbeamertemplate{frametitle}{\vspace{0.55cm}\hspace{0.48cm}\insertframetitle\par\vspace{0.12cm}}
\setbeamerfont{frametitle}{series=\bfseries,size=\LARGE}
\setbeamertemplate{footline}{%
  \begin{tikzpicture}[remember picture,overlay]
    \node[anchor=south,inner sep=0pt] at ([yshift=0.52cm]current page.south)
      {\includegraphics[width=2.64cm,trim=58bp 385bp 58bp 385bp,clip]{assets/lecture/uos-logo-horizontal.png}};
    \node[anchor=east,text=gray,font=\scriptsize]
      at ([xshift=-0.57cm,yshift=0.52cm]current page.south east) {\insertframenumber};
  \end{tikzpicture}\vspace{0.38cm}}
\newcommand{\lab}[1]{\texttt{STM32\_#1}}
\newcommand{\task}[2]{\texttt{STM32\_#1\_과제#2}}
\newcommand{\materialdate}{2026. 10. 06.}
\hypersetup{pdftitle={STM32 4주 실습 계획},pdfauthor={이해리}}
\begin{document}

{\setbeamertemplate{footline}{}
\setbeamercolor{background canvas}{bg=UOSBlue}
\begin{frame}[plain]
  \begin{tikzpicture}[remember picture,overlay]
    \node[anchor=north west,inner sep=0pt,text=white,
      font=\bfseries\fontsize{22}{29}\selectfont,align=left]
      at ([xshift=1.15cm,yshift=-1.5cm]current page.north west)
      {전자전기컴퓨터설계실험Ⅱ\\STM32 4주 실습 계획};
    \node[anchor=north west,inner sep=0pt,text=white,font=\small,align=left]
      at ([xshift=1.15cm,yshift=-4.05cm]current page.north west)
      {2026. 10. 12. -- 11. 02.\\NUCLEO-G474RE / Combo II-DLD};
    \node[anchor=south west,inner sep=0pt] at ([xshift=1.15cm,yshift=0.98cm]current page.south west)
      {\includegraphics[width=1.45cm]{assets/lecture/uos-logo-white.png}};
    \node[anchor=south east,text=white,font=\small\bfseries,align=right]
      at ([xshift=-1.15cm,yshift=0.98cm]current page.south east)
      {작성일 \materialdate\\이해리};
  \end{tikzpicture}
\end{frame}}

% 개념 설명 4쪽: FPGA에서 MCU로, 내부 구조, 프로그램 실행, 제어 방식.
\begin{frame}{FPGA에서 STM32로}
\fontsize{10}{13}\selectfont
\textbf{지난 3주에 구성한 회로를 떠올려 보자.}
\vspace{0.25cm}
\begin{columns}[T]
\begin{column}{0.48\textwidth}
\begin{block}{FPGA}
\fontsize{9}{12}\selectfont
\begin{itemize}\setlength{\itemsep}{0.18cm}
\item Verilog로 회로를 구성한다.
\item 여러 회로가 동시에 동작한다.
\item 카운터·FSM·UART 등을 직접 설계했다.
\end{itemize}
\end{block}
\end{column}
\begin{column}{0.48\textwidth}
\begin{block}{STM32 · MCU}
\fontsize{9}{12}\selectfont
\begin{itemize}\setlength{\itemsep}{0.18cm}
\item 준비된 CPU와 주변장치를 사용한다.
\item CPU가 C에서 변환된 명령을 실행한다.
\item GPIO·타이머·UART·SPI 등을 설정한다.
\end{itemize}
\end{block}
\end{column}
\end{columns}
\vspace{0.30cm}
\textbf{FPGA에서는 회로를 구성했고,\\STM32에서는 준비된 회로를 프로그램으로 제어한다.}
\vspace{0.18cm}
\par\fontsize{8}{11}\selectfont
MCU는 CPU·메모리·주변장치를 하나에 모은 마이크로컨트롤러다.

\end{frame}

\begin{frame}{STM32 안에는 무엇이 있는가}
\begin{center}
\begin{tikzpicture}[scale=0.88,transform shape,
 box/.style={draw=UOSBlue,fill=UOSBlue!6,align=center,font=\footnotesize,
 text width=2.4cm,minimum height=0.82cm},>={Stealth}]
\node[box] (flash) at (-3.5,0) {Flash\\프로그램·초기 데이터};
\node[box,fill=UOSBlue!15] (cpu) at (0,0) {Cortex-M4 CPU\\명령 실행};
\node[box] (sram) at (3.5,0) {SRAM\\변수·스택};
\draw[->,thick] (flash) -- (cpu);
\draw[<->,thick] (cpu) -- (sram);
\node[box,text width=6.4cm] (reg) at (0,-1.5)
 {주변장치 레지스터};
\node[box,text width=8.6cm,minimum height=1.1cm] (hw) at (0,-3.3)
 {주변장치 하드웨어\\GPIO: 디지털 입출력 · 타이머: 시간·PWM\\UART: 비동기 통신 · SPI: 동기 통신\\ADC: 아날로그 입력 · DAC: 아날로그 출력};
\draw[<->,thick] (cpu) -- (reg);
\draw[<->,thick] (reg) -- (hw);
\end{tikzpicture}
\end{center}
\fontsize{8}{10}\selectfont
\textbf{CPU가 주변장치 레지스터를 읽고 쓰면, 하드웨어가 신호를 다룬다.}\\
\texttt{HAL\_GPIO\_WritePin()} → GPIO 출력 레지스터 → GPIO 회로 → 핀의 전압\\
주변장치 레지스터는 CPU 내부의 연산용 레지스터와 다르다.\\
{\scriptsize \href{https://www.st.com/resource/en/datasheet/stm32g474re.pdf}{참고: STM32G474 데이터시트}}
\end{frame}

\begin{frame}{작성한 프로그램이 실행되기까지}
\fontsize{9}{12}\selectfont
\textbf{CPU는 C 문장이 아니라 컴파일된 기계어를 실행한다.}
\begin{center}
\begin{tikzpicture}[scale=0.86,transform shape,
 box/.style={draw=UOSBlue,fill=UOSBlue!6,align=center,font=\footnotesize,
 text width=2.5cm,minimum height=0.85cm},>={Stealth}]
\node[box] (src) at (-3.5,0) {C 소스 작성\\\texttt{.c}};
\node[box] (build) at (0,0) {Build\\컴파일·링크};
\node[box] (elf) at (3.5,0) {ELF 파일\\기계어\\디버깅 정보};
\draw[->,thick] (src) -- (build);
\draw[->,thick] (build) -- (elf);
\node[box] (link) at (3.5,-1.55) {ST-LINK\\SWD로 다운로드};
\node[box] (flash) at (0,-1.55) {STM32 Flash\\프로그램 기록};
\node[box] (start) at (-3.5,-1.55) {리셋 후 실행\\시작 코드 → \texttt{main()}};
\draw[->,thick] (elf) -- (link);
\draw[->,thick] (link) -- (flash);
\draw[->,thick] (flash) -- (start);
\end{tikzpicture}
\end{center}
\fontsize{8.5}{11}\selectfont
\textbf{이번 실습에서는 운영체제 없이 우리 프로그램이 직접 실행된다.}
\begin{itemize}\setlength{\itemsep}{0.08cm}
\item 시작 코드가 실행 환경을 준비하고 \texttt{main()}을 호출한다.
\item 인터럽트가 발생하면 해당 처리 함수가 실행된다.
\item VS Code에서 Build로 실행 파일을 만들고, F5로 다운로드·디버깅한다.
\end{itemize}
\fontsize{7.5}{10}\selectfont
ELF의 프로그램 부분이 Flash에 기록된다. 디버깅 정보는 PC에서 사용한다.
\end{frame}

\begin{frame}{이번 4주에 배울 제어 방식}
\fontsize{8.5}{10.5}\selectfont
\begin{block}{STM32\_0 · 폴링}
\fontsize{8.5}{10.5}\selectfont
CPU가 버튼 상태를 반복해서 확인하고 LED를 제어한다.
\end{block}
\begin{block}{STM32\_1 · 인터럽트·타이머/PWM}
\fontsize{8.5}{10.5}\selectfont
사건이 발생하면 처리한다. 타이머로 일정한 주기의 신호를 만든다.
\end{block}
\begin{block}{STM32\_2 · UART}
\fontsize{8.5}{10.5}\selectfont
PC·STM32·FPGA 사이에서 클록 선 없이 데이터를 주고받는다.
\end{block}
\begin{block}{STM32\_3 · SPI}
\fontsize{8.5}{10.5}\selectfont
클록 신호에 맞춰 STM32·FPGA 사이에서 데이터를 주고받는다.
\end{block}
\vspace{0.02cm}
\textbf{LED·버튼 제어에서 시작해 STM32 ↔ FPGA 통신으로 이어간다.}
\end{frame}

\begin{frame}{4주 일정}
\fontsize{8.5}{11}\selectfont
\textbf{4주 실습에 필요한 매뉴얼은 모두 이 리포지토리에 정리되어 있다.}
\par 매주 별도의 추가 매뉴얼은 제공하지 않으므로,\\목차에서 해당 주차의 문서를 확인하며 진행하면 된다.
\vspace{0.10cm}
\fontsize{8.5}{11}\selectfont
\setlength{\tabcolsep}{3pt}
\renewcommand{\arraystretch}{1.1}
\begin{tabular}{@{}p{12mm}p{17mm}p{34mm}p{35mm}@{}}
\toprule
날짜 & 주차 & STM32 실습 & FPGA 교재 실습\\
\midrule
10/12 & \lab{0} & LED·버튼 폴링\newline 과제1·2 & 곱셈기\newline BCD → 이진수 변환\\
10/19 & \lab{1} & 인터럽트·PWM\newline 과제3·4 & 피아노 → 신호등 → ADC\\
10/26 & \lab{2} & PC·FPGA UART\newline 과제5--7 & 서보 → SRAM → DAC\\
11/2 & \lab{3} & SPI 송신·수신\newline 과제8--12 & SPI 연결 과제에 집중\\
\bottomrule
\end{tabular}
\vspace{0.25cm}
\textbf{\color{UOSBlue}10/26 \lab{2}: 10분 퀴즈를 실시할 수 있음.}\\[0.08cm]
실시 여부와 범위는 별도로 안내한다.
\end{frame}

\begin{frame}[fragile]{실습 프로젝트 받기 · 태그 지정}
\textbf{\color{UOSBlue}배포 예정 태그: \texttt{v1.1.0}}

\small
PowerShell에서 아래 명령을 실행한다. 태그를 지정해\
수업에서 사용하는 버전의 프로젝트를 내려받는다.

\begin{Verbatim}[fontsize=\footnotesize]
git clone --branch v1.1.0 --depth 1 `
    https://github.com/Glaysia/ece2_474.git stm32_lab
code .\stm32_lab\ece2_474.code-workspace
\end{Verbatim}

\begin{itemize}\setlength{\itemsep}{0.12cm}
\item 첫 줄 끝의 \texttt{\char96}는 PowerShell 줄 연결 문자다.
\item \texttt{stm32\_lab}은 새 폴더 이름이다. 이미 있으면 다른 이름을 쓴다.
\item VS Code에서 Debug 선택 → Configure → Build → F5로 실행한다.
\end{itemize}
\vspace{0.10cm}
\footnotesize
\texttt{v1.1.0} 배포 예정. 저장소 공개·태그 푸시 완료 후 위 명령을 사용한다.
\end{frame}
\begin{frame}{10/12 · \lab{0}}
\begin{columns}[T]
\begin{column}{0.56\textwidth}
\textbf{\color{UOSBlue}STM32 · LED와 버튼}
\begin{itemize}\setlength{\itemsep}{0.12cm}
\item 제공 예제: 1초 간격 LED,\\버튼을 누르면 LED 상태 변경
\item \task{0}{1}\\버튼으로 간격 변경:\\1 → 2 → 4 → 8 → 1초
\item \task{0}{2}\\같은 동작을 외부 LED로 구현
\end{itemize}
\end{column}
\begin{column}{0.40\textwidth}
\textbf{\color{UOSBlue}FPGA · 교재 실습 2개}
\begin{enumerate}\setlength{\itemsep}{0.15cm}
\item 4비트 곱셈기
\item BCD → 이진수 변환
\end{enumerate}
\vspace{0.18cm}
\small 각자 장비와 함께 있는 교재를 보고 시뮬레이션과 실제 보드 시연을 수행한다.
\end{column}
\end{columns}
\end{frame}

\begin{frame}{10/19 · \lab{1}}
\begin{columns}[T]
\begin{column}{0.56\textwidth}
\textbf{\color{UOSBlue}STM32 · 인터럽트와 PWM}
\begin{itemize}\setlength{\itemsep}{0.12cm}
\item 버튼 ISR, TIM6 100 ms 타이머,\\PB4 PWM 출력
\item \task{1}{3}\\\task{0}{1}을\\버튼·타이머 ISR로 재현
\item \task{1}{4}\\버튼으로 PWM 비율 변경:\\25\% → 50\% → 75\% → 25\%\\{\small C 디버거: 중단점·한 줄 실행·변수 확인}
\end{itemize}
\end{column}
\begin{column}{0.40\textwidth}
\textbf{\color{UOSBlue}FPGA · 교재 실습 3개}
\begin{enumerate}\setlength{\itemsep}{0.15cm}
\item 피아노
\item 신호등 제어
\item ADC
\end{enumerate}
\vspace{0.18cm}
\small ADC는 FPGA로 외부 변환 장치를 제어한다. STM32 내부 ADC는 사용하지 않는다.
\end{column}
\end{columns}
\end{frame}

\begin{frame}{10/26 · \lab{2}}
\fontsize{10}{13}\selectfont
\begin{columns}[T]
\begin{column}{0.56\textwidth}
\textbf{\color{UOSBlue}STM32 · UART 통신}
\begin{itemize}\setlength{\itemsep}{0.12cm}
\item UART2: PC / UART3: FPGA\\통신 속도는 모두 9600
\item \task{2}{5}\\PC와 MCU 입출력 확인
\item \task{2}{6}\\MCU → FPGA 수신값 확인
\item \task{2}{7}\\UART 명령으로 FPGA LED 제어
\end{itemize}
\end{column}
\begin{column}{0.40\textwidth}
\textbf{\color{UOSBlue}FPGA · 교재 실습 3개}
\begin{enumerate}\setlength{\itemsep}{0.15cm}
\item 서보모터 제어
\item SRAM 읽기·쓰기
\item DAC
\end{enumerate}
\small DAC는 FPGA로 외부 장치를 제어한다.
\end{column}
\end{columns}
\vspace{0.10cm}
\begin{block}{퀴즈 안내}
\small\textbf{10분 퀴즈를 실시할 수 있음}. 여부·범위는 별도 안내.
\end{block}
\end{frame}

\begin{frame}{11/2 · \lab{3}}
\textbf{\color{UOSBlue}SPI · 송신부터 수신 확인까지}
\begin{itemize}\setlength{\itemsep}{0.12cm}
\item \task{3}{8}: MCU 송신 파형을 오실로스코프로 확인
\item \task{3}{9}: MCU → MCU, 수신 측 VCP UART 확인
\item \task{3}{10}: MCU → FPGA, FPGA 수신 회로 작성
\item \task{3}{11}: FPGA → MCU, FPGA 송신 회로 작성
\item \task{3}{12}: FPGA → FPGA, 수신 FPGA UART 확인
\end{itemize}
\vspace{0.10cm}
\small
MCU--FPGA: \textbf{자기 조의 장비 사용} (과제10·11).\\
과제9·12만 조 연결: \textbf{A--B / C--D / E--F / G--H / I--J}.\\
Mode 0 / 8비트 / 높은 비트부터 / SCK 1 MHz 이하.\\
한 방향 연결은 SCK·MOSI·CS·GND 4가닥. I²C는 수행하지 않는다.\\
FPGA 송수신 Verilog는 직접 작성하며, 교재 추가 실습은 없다.
\end{frame}
\end{document}
